\subsection{Extraction parameters, dispersion, and sampling}
\label{ian:hard-sector-smooth}

Status: proved, with no quantum property and no uniform-AA hypothesis: the dispersion trade-off (Corollary~\ref{ian:hard-sector-general-tradeoff}), the bound \(a_0(\varepsilon)\le4\varepsilon^2\) on the dispersion guaranteed by any extraction whose depth depends only on \((\varepsilon,\eta,t)\) and which covers all \(Q:\Omega_n\to[0,1]\) under the uniform law (Proposition~\ref{ian:hard-sector-dispersion-order}), and a dimension-free sampling bound for mixtures of predictors (Remark~\ref{ian:hard-sector-minimax-sampling}). Open: the literal H1 form with one exponent \(C\) in every parameter.

The notation is that of Subsection~\ref{ian:hard-sector-definitions-forward}.
None of the results below is needed for
Theorem~\ref{ian:hard-sector-aa-equivalence}.

\begin{remark}[Why dispersion is obtained]
\label{ian:hard-sector-dispersion-explanation}
The proof of Theorem~\ref{ian:hard-sector-regularity} bins one exactly
calibrated predictor \(g\), and the final approximator is the conditional
mean of \(Q\) on each score bin, not \(g\) itself. Low-variance bins cost at
most \(\sqrt a\), and heavy bins of variance at least \(a\) are bad unless
they are hard; so approximation failure forces a good bin with both mass
at least \(w\) and variance at least \(a\). This handles nearly constant
\(Q\) directly, and the hard sector is a set, so no smooth-to-set rounding
is needed. All conditional means and tests are fixed in unrestricted
preprocessing, and no hereditary simulation hypothesis of the kind used in
Proposition~\ref{tht:reveal-prop} is needed.
\end{remark}

\begin{corollary}[A separately specified dispersion threshold]
\label{ian:hard-sector-general-tradeoff}
Let \(t\ge1\), \(0<\varepsilon<1/8\), \(0<\eta\le1/4\), and
\(0<b\le\varepsilon^2/4\). If
\[
 \ell_\mu\left(
 \left\lceil\frac2{\varepsilon\eta^2b^2}\right\rceil t,Q
 \right)>\varepsilon,
\]
then some set \(H\) satisfies
\[
 \mu(H)\ge\frac{\varepsilon\eta b}{4},
 \qquad
 \operatorname{Hard}(Q,\mu(\cdot\mid H);t,b,\eta b).
\]
Corollary~\ref{ian:hard-sector-explicit} is the case
\(b=\varepsilon^2/4\).
\end{corollary}

\begin{proof}
Apply Theorem~\ref{ian:hard-sector-regularity} with \(a=b\),
\(\gamma=\eta b\), \(M=\lceil1/\gamma\rceil\), \(w=\varepsilon/(2M)\) and
\(K=\lceil2/(\varepsilon\gamma^2)\rceil=\lceil2/(\varepsilon\eta^2b^2)\rceil\).
Then \(0<\gamma<a\le1/4\), \(w<1\), \(\sqrt b\le\varepsilon/2\),
\(Mw/2=\varepsilon/4\) and \(1/(2K\gamma^2)\le\varepsilon/4\), so the
approximation alternative has error at most \(\varepsilon\) and is
excluded. As \(M\le2/\gamma\),
\(w\ge\varepsilon\gamma/4=\varepsilon\eta b/4\).
\end{proof}

\begin{remark}[The literal common-exponent requirement]
\label{ian:hard-sector-common-exponent-warning}
Corollary~\ref{ian:hard-sector-explicit} has independently specified
exponents: depth multiplier \(32\varepsilon^{-5}\eta^{-2}\), inverse
weight at most \(16\varepsilon^{-3}\eta^{-1}\), and dispersion
\(\varepsilon^2/4\). It does not give the literal form
\[
 L=\lceil C(\varepsilon\eta)^{-C}\rceil t,\qquad
 a_0=c\varepsilon^C,\qquad
 \gamma=\eta a_0,
\]
and padding exponents does not reach it: lowering the variance threshold
from \(a\) to \(b<a\) also tightens the covariance tolerance from
\(\eta a\) to \(\eta b\), and \(\operatorname{Hard}\) is not monotone
under this change. For example, under the uniform law on \(\Omega_2\),
fix \(0<\eta\le1/4\) and set
\[
 Q(x)=\frac12+\frac14x_1x_2+\frac{\eta}{16}x_1,
 \qquad a=\frac1{16}.
\]
Then \(Q\in[0,1]\), \(\Var Q=\frac1{16}+\frac{\eta^2}{256}\ge a\), and,
since a depth-one test is an affine function of one coordinate and the
two coordinate covariances are \(\eta/16\) and \(0\),
\(\sup_{F\in\mathcal F_1}|\Cov(Q,F)|=\eta/16=\eta a\).
Hence \(\operatorname{Hard}(Q,\mu;1,a,\eta a)\) holds, but
\(\operatorname{Hard}(Q,\mu;1,b,\eta b)\) fails for every \(0<b<a\).

Corollary~\ref{ian:hard-sector-general-tradeoff} with
\(b=c\varepsilon^C\) gives depth proportional to
\(\varepsilon^{-(2C+1)}\eta^{-2}\) and weight proportional to
\(\varepsilon^{C+1}\eta\), not the common-exponent form.
This obstructs only the padding argument; whether some universal
\(c,C\) exist remains open here.
\end{remark}

\begin{proposition}[The dispersion scale cannot exceed order
\(\varepsilon^2\)]
\label{ian:hard-sector-dispersion-order}
Let \(0<\varepsilon<1/8\), let \(n\ge1\), let \(\mu\) be uniform on
\(\Omega_n\), and put \(Q(x)=\frac12+2\varepsilon\prod_{i=1}^nx_i\).
Then \(\ell_\mu(d,Q)=2\varepsilon>\varepsilon\) for every \(d<n\), while
\(\Var_\nu Q\le4\varepsilon^2\) for every law \(\nu\) on \(\Omega_n\).
Consequently, any extraction statement that applies, for every \(n\),
to the uniform law and to every \(Q:\Omega_n\to[0,1]\) (with no bound
on quantum query cost), and that, from
\(\ell_\mu(L,Q)>\varepsilon\) at a depth \(L\) depending only on
\(\varepsilon,\eta,t\), yields a law \(\nu\) with
\(\Var_\nu Q\ge a_0(\varepsilon)\) (smooth or set, of any weight) must
have \(a_0(\varepsilon)\le4\varepsilon^2\) for all
\(\varepsilon\in(0,1/8)\). In particular, the literal H1 form
with \(a_0=c\varepsilon^C\) is impossible for \(C<2\).
\end{proposition}

\begin{proof}
At every leaf of a tree of depth \(d<n\) some coordinate is unqueried,
so the parity is conditionally uniform on \(\{-1,1\}\); for any leaf
output \(z\),
\[
 \frac12\left|z-\left(\frac12+2\varepsilon\right)\right|
 +\frac12\left|z-\left(\frac12-2\varepsilon\right)\right|
 \ge2\varepsilon,
\]
with equality at \(z=1/2\). Thus \(\ell_\mu(d,Q)=2\varepsilon\). For every
law \(\nu\), \(\Var_\nu Q\le\E_\nu(Q-1/2)^2=4\varepsilon^2\). Applying
the extraction statement with \(n>L\) gives
\(a_0(\varepsilon)\le4\varepsilon^2\); and for fixed \(c>0\), \(C<2\),
small \(\varepsilon\) has \(c\varepsilon^C>4\varepsilon^2\).
The example is quantum-realizable: query all \(n\) coordinates and
accept with probability \(Q(x)\).
\end{proof}

\begin{remark}[What counting does and does not imply]
\label{ian:hard-sector-minimax-sampling}
Uniform empirical control of all Boolean depth-\(t\) test moments by
Hoeffding's inequality and a union bound over the tests of
Lemma~\ref{ian:hard-sector-tree-closure} uses
\(s\ge2\delta^{-2}\log(2N(n,t)/\rho)\) samples for error \(\delta\) and
failure probability \(\rho\), where
\[
 \log N(n,t)=(2^t-1)\log n+2^t\log2.
\]
This argument therefore depends on \(n\), as does the rounding criterion
of Lemma~\ref{ian:hard-sector-rounding}. The loss is real for uniform
empirical test moments: for \(s\) independent uniform points of
\(\Omega_n\), the probability that some coordinate is \(+1\) in every
sample is \(1-(1-2^{-s})^n\), and the empirical mean of that coordinate
then differs from its true mean by one. Simultaneous error below one with
fixed confidence therefore needs \(s=\Omega(\log n)\).

Sampling a mixture of predictors needs no such control. For any
distribution on \(D\in\mathcal D_t\), put \(g=\E D\) pointwise, take
independent \(D_1,\ldots,D_s\) from it, and let
\(\bar D=s^{-1}\sum_iD_i\). Each \(D_i(x)\in[0,1]\) has variance at most
\(1/4\), so the pointwise variance bound and Cauchy--Schwarz give
\[
 \E_{\mathrm{samples}}\E_\mu(\bar D-g)^2\le\frac1{4s},\qquad
 \E_{\mathrm{samples}}\E_\mu|\bar D-Q|
 \le\E_\mu|g-Q|+\frac1{2\sqrt s}.
\]
Some realization achieves this loss, and by
Lemma~\ref{ian:hard-sector-tree-closure} it has depth at most \(st\),
independently of \(n\). So the \(\log n\) count applies only when
simultaneous approximation of all test moments is demanded, not to
sampling a mixture of predictors. This does not by itself give a minimax
proof of sector existence; the extraction of
Theorem~\ref{ian:hard-sector-regularity} uses neither minimax nor a
hardcore lemma from the literature.
\end{remark}
